\section{Contexto do Problema}
% (==>) É assim que esse mecanismo se manifesta concretamente no Brasil e em Mato Grosso.

% Aqui entra muito mais realidade institucional e empírica:
% expansão da soja/pecuária;
% estrutura fundiária;
% Código Florestal;
% fiscalização;
% Amazônia/Cerrado/Pantanal;
% desmatamento;
% precipitação;
% sistemas produtivos;
% instituições brasileiras.

% 2. Contexto econômico e institucional do problema
% 2.1. Produção agropecuária e conversão da vegetação
% Retornos, terra, soja, pecuária, expansão agrícola.

% 2.2. Instituições ambientais e decisões de uso da terra
% Código Florestal, fiscalização, sanções, regularização etc.

% 2.3. Vegetação, clima e produção agropecuária
% Mecanismos biofísicos.

% 2.4. Mato Grosso como ambiente empírico
% Amazônia, Cerrado, Pantanal e heterogeneidade territorial.

Lorem ipsum dolor sit amet, consectetur adipiscing elit. Sed do eiusmod tempor incididunt ut labore et dolore magna aliqua.

Ut enim ad minim veniam, quis nostrud exercitation ullamco laboris nisi ut aliquip ex ea commodo consequat.
